\documentclass[12pt, letterpaper]{article}
\usepackage[margin=1.2in]{geometry}
\usepackage{ebgaramond}
\usepackage[T1]{fontenc}
\usepackage{microtype}
\usepackage{setspace}
\usepackage{parskip}
\usepackage{titling}
\usepackage{fancyhdr}

\setlength{\parindent}{2em}
\onehalfspacing

\pagestyle{fancy}
\fancyhf{}
\fancyhead[L]{\small\textit{The Mirror}}
\fancyhead[R]{\small\textit{NASA Mission Series v1.14}}
\fancyfoot[C]{\small\thepage}
\renewcommand{\headrulewidth}{0.4pt}

\title{\Huge\textbf{The Mirror}\\[0.5em]\large\textit{A short story from Mission 1.14 --- Echo 1}}
\author{NASA Mission Series\\[0.3em]\small Miles Tiberius Foxglove --- Mukilteo, WA}
\date{March 30, 2026}

\begin{document}
\maketitle
\thispagestyle{empty}

\vspace{1em}
\noindent\rule{\textwidth}{0.4pt}
\vspace{1em}

\noindent The evening my father pointed it out to me, I was nine years old and angry about the mosquitoes. We were at the lake --- Conesus, one of the Finger Lakes in western New York --- and the air was the particular August thickness that makes you believe summer will last forever and also that it should stop immediately. The mosquitoes were relentless. I was slapping at my legs and arms and neck and being theatrical about it, the way children are when they want adults to acknowledge that their suffering is unique and historic.

``Come here,'' my father said. He was standing at the end of the dock, looking up. Not at the sky in general --- people look at the sky in general all the time, a vague upward glance that takes in clouds or weather or the general fact of darkness. He was looking at something specific. His head was tilted back and tracking, the way you track a bird, and he had that expression he got when he was about to teach me something. I knew the expression. It meant I was about to learn whether I wanted to or not.

I walked out to the end of the dock, still slapping, and looked where he was pointing.

``There,'' he said. ``See it?''

A star was moving. That was my first thought --- a star had come loose from wherever stars are attached and was sliding across the sky like a marble on a tilted table. It was bright, brighter than any other star visible, a steady silver-white point that moved with the deliberate patience of something that knows exactly where it is going. No blinking. No wavering. Just a smooth, silent arc from roughly west to east, tracking across the summer constellations like it owned them.

``What is it?'' I said, and I had stopped slapping at the mosquitoes.

``Echo,'' my father said. ``It's a satellite. A balloon, actually. A hundred feet across, made of metal-coated plastic, and it's up there right now reflecting radio signals from one side of the country to the other. They launched it twelve days ago.''

\begin{center}
$\bullet\quad\bullet\quad\bullet$
\end{center}

A balloon. I could not reconcile this with what I was seeing. Balloons were red things tied to the backs of chairs at birthday parties. Balloons were the Goodyear blimp drifting over football games, its surface rippling, its motion lazy and wind-dependent. Balloons were not silent, steady, impossibly bright points of light moving faster than any airplane across the meridian of the sky at --- and here my father supplied the number, because he always had the numbers --- eighteen thousand miles per hour.

``It's a thousand miles up,'' he said. ``You're seeing sunlight reflected off the surface. The sun has already set for us, but up there it's still in sunlight. The Mylar surface is like a mirror.''

I did not know what Mylar was. I would learn later that it was a polyester film, drawn thin and coated with vaporized aluminum, so that it reflected ninety-five percent of the light that struck it. I would learn that the balloon had been folded and compressed into a sphere twenty-six inches across, packed inside the nosecone of a Delta rocket, and that when the rocket's payload fairing separated at altitude, gas canisters inflated the balloon to its full one-hundred-foot diameter in about twelve minutes. I would learn that it weighed only a hundred and sixty-seven pounds --- less than my father --- and that its skin was thinner than a human hair. I would learn all of this later. In that moment, standing on the dock at Conesus Lake, what I knew was simpler and larger: there was a mirror in the sky, and I could see it with my own eyes, and it was beautiful.

\begin{center}
$\bullet\quad\bullet\quad\bullet$
\end{center}

My father explained the geometry. Bell Laboratories in Holmdel, New Jersey, had a large horn antenna --- he said it looked like a foghorn the size of a house --- that pointed at the balloon as it passed overhead and transmitted a radio signal at nine hundred and sixty megahertz. The signal traveled upward at the speed of light, struck the balloon's curved surface, and reflected in all directions, the way sunlight scatters off a Christmas ornament. A tiny fraction of that scattered signal continued downward to another antenna at the Jet Propulsion Laboratory in Goldstone, California, twenty-five hundred miles away. The total round-trip time was about eleven milliseconds. In that sliver of time, a human voice could travel from the Atlantic coast to the Pacific coast via a balloon in space.

The word that stuck was \emph{passive}. My father said it with a particular emphasis, the way he emphasized things that mattered. The balloon carried no electronics. No transmitter, no receiver, no amplifier, no computer, no power source. It was a mirror. It did nothing but exist in the right shape at the right altitude and reflect what was sent to it. The intelligence was on the ground --- the massive antennas, the sensitive receivers, the powerful transmitters. The balloon's only job was to be there and be shiny.

I thought about this for a long time afterward. There was something clean about it, something I could not articulate at nine but that I now understand as elegance. A system in which the space segment does nothing --- literally nothing, not even maintaining its own shape against the vacuum --- and all the work happens at the endpoints. The balloon was infrastructure in the purest sense: a dumb, beautiful, reflective surface that enabled communication without participating in it. A bridge that carried no weight of its own. A window that generated no light.

\begin{center}
$\bullet\quad\bullet\quad\bullet$
\end{center}

We watched it for perhaps four minutes. It crossed the sky from the constellation Lyra to somewhere in the direction of Pegasus --- I was not good with constellations then --- and then it faded. Not suddenly, like a light switching off, but gradually, the way a sunset fades: the brightness diminishing over several seconds until the point of light became indistinguishable from the background stars and then disappeared entirely. It had entered Earth's shadow. The mirror was still there, still in orbit, still spinning slowly on its axis, but the sunlight that made it visible had been blocked by the planet itself.

``It's still there,'' my father said, reading my face. ``You just can't see it. It'll come back next orbit.''

Something about a Portuguese man o' war --- I had seen pictures in a \emph{National Geographic} at the dentist's office. The float, the pneumatophore, was a gas-filled bladder that rode the surface of the ocean and caught the wind like a sail. It had no engine, no rudder, no navigation. It went where the wind and current took it, trailing its venomous tentacles below. It was, in a sense, a passive system: a translucent balloon on the interface between two media, responding to forces it could not control, catching what the environment offered. Echo 1 was a metallic balloon on the interface between atmosphere and vacuum, catching radio waves it could not process. Both were floats. Both were mirrors of a kind --- the man o' war reflecting and refracting light through its gas-filled body, the satellite reflecting radio waves off its metallic skin. Both were carried by currents they did not resist.

I did not think about the man o' war that night. I thought about it years later, when I was studying physics and encountered the wave equation in Schr\"odinger's formulation. The wave function, $\psi$, describes a quantum system as a superposition of states --- the particle exists everywhere and nowhere until observed. Until you collapse the wave function by measuring it, the system is pure possibility, a smear of probability density across all of space. Erwin Schr\"odinger, born in Vienna in 1887, won the Nobel Prize in 1933 for this insight and spent the rest of his life being uncomfortable with it. His famous cat was a protest: he did not believe that a cat could be simultaneously alive and dead. He thought the wave function was a mathematical convenience, not a physical reality.

\begin{center}
$\bullet\quad\bullet\quad\bullet$
\end{center}

But standing on the dock, watching that bright point slide across the sky, I was already doing the measurement. The photons that entered my eyes had left the Sun eight minutes earlier, traveled ninety-three million miles to Echo 1, reflected off the Mylar surface, and traveled another thousand miles down to the retina of a nine-year-old standing on a dock in western New York. The wave function of those photons --- their quantum-mechanical description as superpositions of all possible paths --- collapsed when they hit my retinal cells and triggered a nerve impulse. Before I saw the satellite, it existed as a probability cloud. After I saw it, it was real, located, definite: \emph{there}, moving from Lyra toward Pegasus, at an elevation of about sixty degrees, with an apparent magnitude of roughly negative one. My observation made it concrete.

Schr\"odinger would have hated this interpretation. He would have said that the satellite was obviously there whether I looked at it or not, that the wave function is a tool for calculation and not a description of reality, that the Moon does not disappear when nobody is watching. And he would have been right, in the pragmatic sense that lets us build bridges and launch satellites. But there is a deeper sense in which the observation matters: the satellite had no meaning until someone saw it. It was a mirror, and a mirror without an observer is just a surface. The reflection has to go somewhere. The signal has to be received. The wave function has to collapse into a measurement. The whole point of Echo 1 was not to be in orbit --- any rock can orbit --- but to be seen, to be used, to reflect a signal that someone on the ground was waiting for. Without the receiver, the transmitter is just shouting into empty air. Without the observer, the mirror reflects nothing that matters.

My father and I stood on that dock and we were the receivers. The ground stations at Holmdel and Goldstone were the receivers of the radio signal, but we were the receivers of the light, and in that moment the system was complete: Sun to satellite to eye. Transmitter, mirror, receiver. Source, medium, detector. The wave function, the reflection, the measurement.

\begin{center}
$\bullet\quad\bullet\quad\bullet$
\end{center}

I went back to the dock every clear night for the rest of that August. I learned to predict when Echo 1 would be visible --- my father showed me how to use the orbital period (about two hours and eighteen minutes) and the geometry of sunlight and shadow to estimate the passes. Some nights it appeared in the west and crossed the entire sky, taking five minutes to traverse my visible horizon. Other nights it appeared briefly, climbing only a few degrees above the trees before entering shadow. Once, on a night with no moon and exceptional transparency, I saw it flare --- a brief, sharp brightening as the satellite's orientation brought a flat section of Mylar perpendicular to my line of sight, creating a momentary specular reflection that was brighter than Venus. It lasted perhaps half a second, a flash of silver-white fire in the constellation Cygnus, and then it dimmed back to its usual steady glow.

The mosquitoes were still there. I no longer cared about them.

I learned later that thousands of people watched Echo 1 that summer. It was, for a few months, the most visible artificial object in the night sky --- brighter than any star, moving with a stately certainty that distinguished it from the frantic blinking of aircraft. Newspapers published its schedule. Families went outside after dinner to watch it pass. For millions of people, Echo 1 was the first satellite they ever saw with their own eyes. It made space real in a way that Sputnik's beeping radio signal had not. You could \emph{see} it. You could point at it and say: we put that there.

A mirror in the sky. A balloon that weighed less than a man. A surface thinner than a hair, reflecting the light of the Sun and the voices of the Earth across a thousand miles of vacuum. No computer, no engine, no intelligence of any kind. Just shape, and surface, and the patient physics of reflection.

\vspace{1.5em}
\noindent\textit{The man o' war floats. The wave function evolves. The mirror reflects. Somewhere between the signal and its echo, between the question and its answer, between the observation and the thing observed, there is a moment when the universe acknowledges that you are watching --- and reflects your gaze back to you, eleven milliseconds later, from a balloon made of aluminum and plastic and the audacity to put a mirror in the sky.}

\vspace{2em}
\begin{center}
\rule{3cm}{0.4pt}\\[1em]
\textit{For Erwin Schr\"odinger (1887--1961), who knew the cat was always alive}\\[0.5em]
\small NASA Mission Series v1.14 --- Echo 1\\
Miles Tiberius Foxglove --- Mukilteo, WA\\
March 30, 2026
\end{center}

\end{document}
